%======================================================================
% USAGE :
%   \begin{boxDefault}[Titre optionnel] ... \end{boxDefault}   % boîte grise sans icône
%   \begin{boxInfo}[Titre optionnel] ... \end{boxInfo}
%   \begin{boxSuccess}[Bravo !] ... \end{boxSuccess}
%   \begin{boxWarning}[Attention] ... \end{boxWarning}
%   \begin{boxDanger}[Important] ... \end{boxDanger}
%   \begin{boxTeacher}[Note pour l'enseignant] ... \end{boxTeacher}
%   \begin{boxDownload}[Ressources] ... \end{boxDownload}
%   \begin{boxHelp}[Indice] ... \end{boxHelp}
%   \begin{boxObjectif}[Objectifs] ... \end{boxObjectif}
%   \begin{boxValidation}[Critères] ... \end{boxValidation}
%======================================================================

%----------------------------------------------------------------------
% Gestion du chemin des icônes (\DocDisplayBoxPictures)
%----------------------------------------------------------------------
% Si DocSetup n'a pas défini \DocDisplayBoxPictures, on met un défaut.
% Si une ancienne macro \displayboxpictures existe, on l'utilise comme source.
\makeatletter
\@ifundefined{DocDisplayBoxPictures}{%
	\@ifundefined{displayboxpictures}{%
		% ni DocDisplayBoxPictures ni displayboxpictures -> défaut
		\providecommand{\DocDisplayBoxPictures}{../_includes/boxes_images}%
	}{%
		% ancienne macro displayboxpictures définie -> on la réutilise
		\let\DocDisplayBoxPictures\displayboxpictures%
	}%
}{}%
\makeatother

%----------------------------------------------------------------------
% Packages nécessaires (tolérant si déjà chargés ailleurs)
%----------------------------------------------------------------------
\RequirePackage{graphicx}
\RequirePackage{marginnote}
\RequirePackage{tcolorbox}
\RequirePackage{tikz}
\RequirePackage{verbatim}
\RequirePackage{etoolbox}

% Librairies tcolorbox utiles
\tcbuselibrary{skins,breakable}

%======================================================================
% Macro utilitaire pour déclarer les boîtes standard
%======================================================================
% \CFPTDeclareBox{<NomEnv>}{<prefixCouleur>}{<colTitreTexte>}{<épaisseurCadre>}{<icone>}
%
% Paramètres :
%  - <NomEnv>        : nom de l'environnement tcolorbox à créer
%                      (boxInfo, boxSuccess, boxWarning, …).
%  - <prefixCouleur> : préfixe des couleurs définies dans inc_colors :
%                      <prefix>Background, <prefix>Color, <prefix>TitleBackground.
%                      Ex. "info", "success", "warning", "danger", "teacher",
%                      "download", "aide", "objectif", "validation", "default".
%  - <colTitreTexte> : couleur du texte du titre (ex. white, black).
%  - <épaisseurCadre>: épaisseur du cadre (ex. 0.5pt, 1.0pt).
%  - <icone>         : nom de fichier icône (ex. box-info.png) dans
%                      \DocDisplayBoxPictures. Chaîne vide ("") = aucune icône.
%
% Comportement spécifique :
%  - Les boîtes sont "breakable" (coupure possible sur plusieurs pages).
%  - L'icône est dessinée en overlay en haut à droite :
%      * une fois pour une boîte non cassée (overlay unbroken),
%      * une fois sur le premier segment pour une boîte cassée (overlay first).
%  - La marge droite (clé 'right') dépend de la présence d'icône et de titre :
%      * pas d'icône (#5 vide)              -> right = 0.5cm
%      * icône + titre présent              -> right = 0.5cm
%      * icône mais PAS de titre            -> right = 1.0cm
%    -> le texte ne passe pas sous l'icône lorsqu'il n'y a pas de titre.
%======================================================================

\makeatletter
\newcommand{\CFPTDeclareBox}[5]{%
	\newtcolorbox{#1}[1][]{%
		colback=#2Background,
		colframe=#2Color,
		colbacktitle=#2TitleBackground,
		coltitle=#3,
		boxrule=#4,
		width=\linewidth,
		enhanced,
		breakable,
		%--------------------------------------------
		% Marge à droite :
		%  - pas d'icône (#5 vide)           -> 0mm
		%  - icône + titre présent          -> 0mm
		%  - icône mais PAS de titre        -> 1.2cm
		%--------------------------------------------
		right=\ifx&#5&%
					0.5cm%
				\else
					\iftcb@hasTitle%
						0.5cm%
					\else
						1.0cm%
					\fi
				\fi,
		title={##1},
		fonttitle=\bfseries,
		toptitle=2mm,
		bottomtitle=2mm,
		% cacher l'entête si pas de titre
		title code={\if\relax\detokenize{##1}\relax\else\tcbtitle\par\fi},
		% Icône : un seul dessin (boîte non cassée ou 1er segment)
		overlay unbroken={%
			\ifx&#5&%
			\else
			\node[anchor=north east] at (frame.north east)
			{\includegraphics[height=0.70cm]{\DocDisplayBoxPictures/#5}};%
			\fi
		},
		overlay first={%
			\ifx&#5&%
			\else
			\node[anchor=north east] at (frame.north east)
			{\includegraphics[height=0.70cm]{\DocDisplayBoxPictures/#5}};%
			\fi
		},
	}%
}
\makeatother

%======================================================================
% Déclaration des boîtes standard (couleurs centralisées dans inc_colors)
%======================================================================

% boîte par défaut grise, sans icône
\CFPTDeclareBox{boxDefault}{default}{white}{0.5pt}{}

% info
\CFPTDeclareBox{boxInfo}{info}{white}{0.5pt}{box-info.png}

% success
\CFPTDeclareBox{boxSuccess}{success}{white}{0.5pt}{box-success.png}

% warning (titre en noir pour lisibilité sur fond jaune)
\CFPTDeclareBox{boxWarning}{warning}{black}{0.5pt}{box-warning.png}

% danger
\CFPTDeclareBox{boxDanger}{danger}{white}{0.5pt}{box-danger.png}

% teacher
\CFPTDeclareBox{boxTeacher}{teacher}{white}{0.5pt}{box-teacher.png}

% download (cadre plus épais)
\CFPTDeclareBox{boxDownload}{download}{white}{1.0pt}{box-download.png}

% aide
\CFPTDeclareBox{boxHelp}{aide}{white}{0.5pt}{box-aide.png}

% Objectif
\CFPTDeclareBox{boxObjectif}{objectif}{black}{0.5pt}{box-objectif.png}

% Validation
\CFPTDeclareBox{boxValidation}{validation}{black}{0.5pt}{box-validation.png}